\section{Method}
\label{sec:method}

Given a screenshot $I$ and an instruction $q$, we produce an ordered list of options
$(c_1,\dots,c_K)$, each a box with a role (text or icon) and, for the options kept, a label. A separate
\emph{selector} model receives the list and picks one option. We do not train or modify the selector.
The pipeline has three stages: candidate generation (\cref{sec:candidates}), instruction-aware ranking
(\cref{sec:ranker}) and labelling of the kept icons.

\subsection{Candidates from pixels}
\label{sec:candidates}

\paragraph{Sources.}
\emph{Text:} OCR lines (Apple Vision, accurate mode) give boxes and their strings for free.
\emph{Elements:} a detector we fine-tune on accessibility labels (below), run on the full screenshot plus
overlapping tiles of about $1280$ px when the screen is large, so that 16 px icons on a 4K screen keep their
detail. Tile detections that touch a tile edge are dropped, and the rest are merged by NMS.
\emph{Icons:} OmniParser v2's icon detector~\cite{lu2024omniparser}, run the same way, with its confidences
halved so that our detector's boxes take precedence where both fire.

\paragraph{A detector trained on free accessibility labels.}
When an application exposes its AX tree, the tree is a free but incomplete source of labels. We collected
19{,}898 deduplicated macOS screenshots of 115 applications and websites with their AX trees and split them
\emph{by application} (80 / 12 / 23 apps for train / val / test). The labels are the visible actionable AX
elements. To fill in controls that AX omits, we add boxes from a public GUI detector (GPA-GUI-Detector) with
confidence $\ge 0.5$ that overlap no AX box (IoU $<0.3$). We leave out screens with almost no AX elements, which
would teach the detector to ignore real controls. We fine-tune GPA-GUI-Detector (YOLO11m) on these labels at
$1280$ px for 15 epochs. On the 3{,}275 screens of the 23 held-out applications (\cref{tab:detector}), the
fine-tuned detector boxes more of the AX elements than the public detectors (72.6\% at IoU${\ge}0.5$ vs.\ at most
61.6\%) and is the fastest. Training on AX boxes alone gives the best box quality. Adding the GPA boxes raises the
share of AX elements whose centre falls inside some detection, which is what a click needs, from 90.0\% to 97.3\%,
at 65 instead of 49 boxes per screen.

\begin{table}[t]
\centering
\small
\setlength{\tabcolsep}{3.5pt}
\begin{tabular}{l rrrrr}
\toprule
Detector (1280 px) & mAP50 & R@.5 & centre & boxes & ms \\
\midrule
GPA-GUI-Detector & 37.1 & 56.2 & 95.7 & 58 & 57 \\
OmniParser v2 & 52.1 & 61.6 & 78.9 & 37 & 42 \\
Screen2AX$^\ast$ & 17.8 & 40.9 & 71.6 & 59 & 34 \\
Ours, AX labels only & \textbf{66.7} & 71.1 & 90.0 & 49 & 21 \\
Ours, AX + GPA labels & 62.1 & \textbf{72.6} & \textbf{97.3} & 65 & 21 \\
\bottomrule
\end{tabular}
\caption{\textbf{Detectors on held-out applications} (3{,}275 screens, 142k AX elements). R@.5: AX boxes matched at
IoU${\ge}0.5$; centre: AX boxes whose centre lies inside a detection; boxes: median per screen; ms: per image on an
A10G, untiled. $^\ast$At 640 px, its better setting.}
\label{tab:detector}
\end{table}

\paragraph{Merging.}
Sources are merged into one list. OCR lines enter first. A detection whose IoU with an existing item exceeds $0.3$,
or that encloses the centres of items no more than $1.5\times$ its area, only adds its source tag and confidence to
them; otherwise it becomes a new \emph{icon} item. Text items carry their OCR string; icons have no label yet.

\subsection{Instruction-aware ranking}
\label{sec:ranker}

A screen of a professional application yields 300--600 candidates (\cref{tab:main}). Captioning every icon is the
most expensive step of pixel parsing, and a long list also makes the selector's choice harder. We therefore rank
all candidates against the instruction first and keep only the top $K$.

\paragraph{Scoring.}
Each candidate $c$ is shown to an image encoder $f_I$ as a \emph{context crop}: a square window centred on the box,
three times its longer side (at least 96 px), clipped to the screen and resized to the encoder's input size, with
the candidate outlined in red (line width about $1/112$ of the window). The outline tells the model which element
is meant~\cite{shtedritski2023redcircle}, and the window shows its neighbours, \eg the label next to an icon or the
toolbar it sits in. The instruction is encoded by the text encoder $f_T$ and
\begin{equation}
  s(q,c) = \tau\,\langle f_T(q), f_I(\mathrm{crop}(I,c)) \rangle ,
\end{equation}
with unit-normalised embeddings and SigLIP's learned temperature $\tau$. Image embeddings are computed once per
screen and shared by every instruction on it; each new instruction costs one text embedding and a matrix product.

\paragraph{Per-screen listwise training.}
A list is good when the target ranks high among the \emph{other candidates of the same screen}, so we train
for that rank directly. For an instruction $q$ on a screen with candidates $C$ and matched targets $T\subseteq C$,
\begin{equation}
  \mathcal{L}(q) = -\log \sum_{t\in T} \frac{\exp s(q,t)}{\sum_{c\in C} \exp s(q,c)} ,
\end{equation}
which counts any matched target as correct (a target box can match several overlapping candidates). All other
boxes on the screen, including near-identical neighbouring icons, act as hard negatives. We compare against
SigLIP's pointwise sigmoid loss over the same pairs in \cref{sec:ablations}.

\paragraph{Training data.}
Instructions and target boxes come from public grounding sets that are disjoint from our evaluation benchmarks:
OS-Atlas desktop (macOS and Linux)~\cite{wu2024osatlas}, MMBench-GUI~\cite{wang2025mmbenchgui},
OmniACT~\cite{kapoor2024omniact} and WebLINX~\cite{lu2024weblinx}. We run our candidate generator on their
screenshots and mark as targets the candidates whose centre falls in the answer box or that overlap it with
IoU $>0.3$ (and are at most $4\times$ its area). This gives 4{,}816 training screens with 63{,}141 instructions;
253 further screens are held out to select the checkpoint by recall@10. Per step we keep every target and sample
up to 96 candidates and 16 instructions per screen.
We fine-tune SigLIP-Base/16 at $384$ px~\cite{zhai2023siglip} with AdamW (lr $10^{-5}$, weight decay $0.05$,
one-cycle schedule, bf16) for three epochs; held-out recall@10 rises from $0.749$ (zero-shot) to $0.958$.

\paragraph{Labelling.}
Text options are labelled with their OCR string. Only icons that survive the cut are captioned, with OmniParser's
Florence-2 captioner on the icon crop, so the captioning cost grows with $K$ and not with the number of elements
on the screen. Batched on an A10G, captioning costs about 2~ms per icon.
